\textbf{What supports the proposed limitation.} Holding marked inputs fixed and changing their reconstruction targets improved digit clustering in every seed. Adding the geometry hurt DEC as well as IDEC. Some of the performance loss was already present at their shared initialization. Keeping the reconstruction term during clustering was associated with a modest additional mean loss in this comparison, with a reversal in one seed. The experiment supports a contribution from the reconstruction objective while leaving the fraction of the accuracy loss attributable to background encoding unresolved.

\textbf{Where to apply the spatial prior.} Among the three Gaussian placements, weighting only pretraining gave the highest marked accuracy (68.80\%). Its marked gain of $3.66\pm4.74$ points over IDEC was positive in four seeds, with a mean clean gain of 0.84 points. Weighting only fine-tuning reduced marked ACC by $2.26\pm1.29$ points even though its decoder response also fell, to 0.00480. For centered digits, the pretraining-only variant offers the best observed tradeoff among the weighted objectives. This is a stage-ablation finding on the same split. The prespecified two-stage variant had a small, inconsistent gain in marked accuracy.

\textbf{Skipping reconstruction entirely.} With a random encoder, KL-only clustering scored 31.00\% clean and 28.90\% marked. In seed 4 both domains collapsed to one cluster (11.35\% ACC). A 2.10-point mean drop from this low baseline gives little evidence of useful robustness. The standard DEC comparison in Table~3 uses a trained encoder and retains all ten clusters.
